% !PRINTED 686
\weberpage{686}
It is now to be proved that $\mathfrak{L}$ is divisible by the branch polygon $\mathfrak{Z}_\beta$.

Let us first assume that $\mathfrak{A}$ and $\mathfrak{B}$ are relatively prime to $\mathfrak{Z}_\beta$; then there is an integral rational function $x$ of $\beta$ which vanishes at no point of $\mathfrak{Z}_\beta$ and for which $\omega = x\alpha$ becomes an integral function of $\beta$. If we set
\[
f(\omega) = x^n F(\alpha,\beta),
\]
then:
\[
f'(\omega) = x^{n-1}F'(\alpha).
\]

If $\tau$ is the basis form of $\Omega$ with respect to $\beta$, then by \S\ 182 $f'(\omega)$ is divisible by $f'(\tau)$; and since by assumption $\beta$ becomes infinite at no branch point, $f'(\tau)$, and consequently $F'(\alpha)$ also, is divisible by the branch polygon $\mathfrak{Z}_\beta$, and our assertion is proved. We set:
\[
\mathfrak{L} = \mathfrak{R}\mathfrak{Z}_\beta,
\]
\weq{8}{F'(\alpha) = \dfrac{\mathfrak{R}\mathfrak{Z}_\beta}{\mathfrak{A}^{n-2}\mathfrak{B}^m}.}

This formula is thereby proved only under the assumption that neither $\alpha$ nor $\beta$ becomes infinite at one of the branch points.

But if we make the linear substitution:
\weq{9}{\begin{aligned}
\alpha(\alpha_1+1) &= a\alpha_1,\\
\beta(\beta_1+1) &= b\beta_1,
\end{aligned}}
then by \S\ 190
\[
\mathfrak{Z}_\beta = \mathfrak{Z}_{\beta_1},
\]
and we can choose the constants $a$, $b$ so that $\alpha_1$ and $\beta_1$ become infinite at no point of $\mathfrak{Z}_\beta$. We have only to take for $a$ and $b$ any constant values that $\alpha$ and $\beta$ do not assume at any point of $\mathfrak{Z}_\beta$.

If $F_1(\alpha_1,\beta_1)=0$ is the rational equation subsisting between $\alpha_1$, $\beta_1$, then:
\[
F_1(\alpha_1,\beta_1) = (\alpha_1+1)^n(\beta_1+1)^m F(\alpha,\beta),
\]
or also, by (9):
\[
\alpha^n\beta^m F_1(\alpha_1,\beta_1) = a^n b^m \alpha_1^n \beta_1^m F(\alpha,\beta);
\]
further $d\alpha : d\alpha_1 = \alpha^2 : a\alpha_1^2$, and consequently, having regard to $F=0$, $F_1=0$:
\weq{10}{\alpha^{n-2}\beta^m F_1'(\alpha_1) = a^{n-1}b^m \alpha_1^{n-2}\beta_1^m F'(\alpha).}
